\begin{afterword}
\summaryhead

The post starts from a joke question: how can you be seen as generous without spending much? Hsee found that a \$45 scarf is seen as a more generous gift than a \$55 coat. The author presents this as a case of a general effect. In Hsee's dictionary study, a like-new dictionary with 10,000 entries was valued above a torn one with 20,000 entries when each was seen alone, and below it when the two were seen side by side. A number of entries means little on its own, while a torn cover is plainly bad.

Slovic and colleagues found the same with gambles. A 7/36 chance to win \$9 became more attractive when a 29/36 chance to lose 5 cents was added, because the small loss gives the \$9 a scale. In Hsee's ice cream study, an overfilled small cup beat a larger, underfilled one when each was seen alone. The post ends with tongue-in-cheek advice: to look generous on a fixed budget, buy an expensive item from a cheap class of goods.

\responsehead

The post reports its studies accurately. The dictionary, gamble and ice cream figures match their sources, and the effect it describes, that options judged one at a time are judged by whatever attribute is easy to evaluate, has been replicated: a 2023 replication of Hsee's gift, ice cream and dinnerware studies found the ``less is better'' effect in all three, though more weakly for the ice cream.

The weakest example is the one the post opens with. Hsee's gift study had no side-by-side condition, so it did not show the reversal the post calls its general phenomenon. When the 2023 replication added one, most people shown both gifts still chose the scarf, and the replicators allow that the cheapest coat may simply be seen as a poor coat. The post's sentence also leaves out what made the scarf look generous: respondents were told it was near the top of a \$5 to \$50 range, and the coat near the bottom of a \$50 to \$500 range. The post gives that principle only near its end.

One statement is firmer than its source: ``it only works if the subjects don't see the two gambles side-by-side'' is a reasonable prediction that the gamble studies did not test.

The shopping advice is a joke, and its core is Hsee's own conclusion. It depends on the recipient knowing roughly what the class of goods costs, which Hsee's respondents were told and which the post assumes.

In short: an accurate report of a finding that has replicated, opened with a gift study that did not test the reversal it is used to illustrate.
\end{afterword}
